\section{Introduction}
\label{sec:introduction}

Geostationary satellites observe the Indian monsoon every 15--30 minutes and are the only source of continuous cloud and rainfall information over much of South Asia. Their images are geolocated as if every pixel lay on the Earth's surface. A cloud top several kilometres high lies on the same line of sight but closer to the satellite, so it appears displaced away from the sub-satellite point by about $h\tan(\mathrm{VZA})$, where $h$ is the cloud-top height and VZA the viewing zenith angle. This parallax displacement is zero below the satellite and grows towards the edge of the disk.

India is unusually exposed to this effect. INSAT-3DR, the operational Indian geostationary imager, sits at 74\,\textdegree E, close to the western edge of the country. Its viewing angle reaches 33--38\textdegree{} over the Himalaya and Northeast India, where monsoon convection reaches the tropopause at 15--17\,km. Deep convective tops there appear 8--12\,km from their true position, two to three infrared pixels. Products that use INSAT-3DR imagery, from rainfall estimates to cloud masks and machine-learning classifications of precipitation type, inherit this displacement, and comparisons with radar or rain gauges mix it with any real error.

Parallax correction is well established elsewhere. \citet{vicente2002} showed its importance for satellite rainfall estimates, and their correction has been used operationally in the GOES Hydro-Estimator since 2001 \citep{scofield2003}. \citet{bielinski2020} derived an accurate geometric solution on the ellipsoid and validated it against ground radar over Europe. An open-source correction exists in satpy \citep{raspaud2018,satpy_parallax}, which assumes a spherical Earth and does not handle terrain. For INSAT-3D and INSAT-3DR, however, we found no published parallax correction or validation.

Validating a parallax correction is harder than it seems. A radar measures precipitation, not cloud top, and geolocation errors of the imager mix with the parallax. ISRO's data-quality reports measure INSAT-3DR geolocation error against ground control points; for full-disk images in June 2020 it was 1.7--2.7 infrared pixels to the north \citep{mosdac_dqe2020}, comparable to the parallax itself. How the INSAT--radar offset responds to a parallax correction, and how much of the correction a radar can verify, has not been examined for India.

This paper addresses both the correction and its verification. Our contributions are:
\begin{enumerate}
  \item \textbf{An open-source tool, PyGeoParallax-IN} \citep{pygeoparallax}. It corrects INSAT-3DR imagery for parallax by exact ray tracing on the WGS84 ellipsoid, optionally with terrain, and applies a navigation correction measured from coastlines. It produces uncorrected and corrected datasets on the same grid with occlusion-aware resampling.
  \item \textbf{A validation against GPM Ku-band radar} in 63 coincident monsoon cases (2018--2024; 38 usable) across four Indian regions that span the full range of viewing angles. It uses two independent offset measures, at the scale of the storm and of the convective core.
  \item \textbf{A test of what radar validation can resolve.} Most of the computed parallax is not visible against GPM at convective cores. With a synthetic test and a navigation correction fitted to GPM, we show that this is not a method error and cannot be explained by navigation alone, and that storm structure limits GPM-based validation of absolute position to about $\pm$3--5\,km.
  \item \textbf{A characterisation of INSAT-3DR L1C navigation error} and of the remaining error sources: cloud-top height (standard atmosphere versus ERA5), terrain (SRTM) and the slanted GPM radar beam.
\end{enumerate}

Section~\ref{sec:data} describes the data, Section~\ref{sec:method} the correction method and Section~\ref{sec:validation} the validation design. Section~\ref{sec:results} presents the results, Section~\ref{sec:discussion} discusses them, and Section~\ref{sec:conclusions} concludes.
